Pre-trained language models encode knowledge through learned token embeddings that map surface forms to semantic vectors. We investigate what happens to a transformer's internal representations when this mapping is systematically disrupted by a bijective permutation of the vocabulary at inference time. Using GPT-2 Small (124M parameters) and WikiText-2, we apply a fixed random bijection $\perm: \vocab \to \vocab$ to all token IDs and compare residual stream activations against normal processing across all 12 layers. We find that the model \emph{never} recovers from embedding misalignment. The residual stream exhibits a distinctive U-shaped cosine similarity curve---moderate similarity at the embedding layer (0.59), dropping to a minimum at layers 6--7 (0.47), then rising to 0.85 at the final layer. However, this late-layer convergence is not semantic recovery: logit lens analysis shows near-zero prediction match ($<$1\%) at all layers, and perplexity increases 2,315$\times$. The convergence instead reflects structural regularities in GPT-2's output formatting. Furthermore, similarity \emph{decreases} with more context (Spearman $\rho = -0.98$, $p < 0.001$), contradicting the hypothesis that in-context processing might decode the cipher. Our results empirically ground the theoretical impossibility of permutation inversion in causal transformers and reveal a two-phase structure in the residual stream: identity-dependent processing in layers 0--9 and identity-independent output formatting in layers 10--11.
